
\section{Homology and the fundamental group}\label{chapter:stallings}
Adams' theorem allows one to use homology and the cobar construction to approximate the fundamental group:
\begin{thm}[\cite{Stallings1975}, Theorem 5.4]\label{thm:stallings}
    Let $R$ be a ring and $X$ a path-connected space with single 0-simplex $x$. Let $\mathcal{I}$ be the augmentation ideal of the group ring $R[\pi_1(X,x)]$. Let $E$ be the spectral sequence associated to the one-sided cobar bicomplex on $C_\bullet(X)$ with basepoint $x$, and the filtration by tensor length. Then there is an isomorphism
    \[E^\infty_{-p,p} \to \mathcal{I}^p/\mathcal{I}^{p+1}\]
\end{thm}
A dual version of this theorem with cohomology and the bar construction can be found in \cite{gadish_letterbraiding}, Theorem 1.2(3). The goal of this section will be to prove the following integral version of this theorem for the two-sided cobar construction:
\begin{thm}\label{thm:stallings_generalization}
    Let $X$ be a path-connected space with basepoints $a,b$. Let $\mathcal{I}$ be the augmentation ideal of the group ring $\Z\pi_1(X,a)$. Let $E$ be the spectral sequence associated to the two-sided cobar construction 
    \[\Cobar(b,X,a):= \Cobar(N_\bullet(\{b\}), N_\bullet(X), N_\bullet(\{a\}))\]
    where the filtration is given by tensor length. Then there is an isomorphism
    \[E^\infty_{-p,p}\to \Z\pi_X(a,b)\mathcal{I}^p/\Z\pi_X(a,b)\mathcal{I}^{p+1}.\]
\end{thm}
More precisely, we will construct a collection of maps 
\[\Z\pi_X(a,b)/\Z\pi_X(a,b)\mathcal{I}^{p+1} \to H_0(\Cobar^{\leq n}(b,X,a))\]
for any $a,b \in X^0$ which assemble into a functor between two suitably defined categories whose objects are the zero-simplicies of $X$, and morphisms encode the cobar construction / fundamental groupoid ring data. 

\subsection{Setup}
We will imitate Stallings' strategy. He relies on the following result of Brown:
\begin{thm}[Theorem 10.2, \cite{Brown1959}]
    Let $C$ be a dga $R$-coalgebra. Suppose $C_0  = R$ and $H_1(C) = 0$. Then there is a chain equivalence
    \[\Cobar(\Barr(C)) \to C.\]
\end{thm}
 This is an intermediate step and we don't discuss the bar construction again, so we refer the reader to sources such as \cite{Milgram1967}, Chapter 16 of \cite{Brown1982}, or Chapter 1.5 of \cite{May1999-tb}. In our case we specialize to the case where $C$ is the acyclic cochain complex consisting of the integral group ring $\Z G$ concentrated in degree zero. This is useful for the following reason.
\begin{prop}\label{prop:cobar_zg}
    There is a chain equivalence 
    \[\Cobar(C_\bullet(BG)) \to \Z G\]
    and consequently for any path-connected space $X$ with a single 0-cell $a$, there is an isomorphism 
    \[H_0\Cobar(C_\bullet(X)) \to \Z\pi_1(X,a).\]
    \begin{proof}
        The bar construction on $\Z G$ is equivalent to singular chains on the classifying space $BG$. Then apply Brown's theorem.
        The second isomorphism follows by Corollary 5.3 of \cite{Stallings1975}, which says that the map of $X$ into the $K(\pi_1(X,a), 1)$ induces an isomorphism on cobar constructions.
    \end{proof}
\end{prop}
To be specific, Stallings shows that the map 
\[\tot_0(\Cobar(C_\bullet(X))) \to \Z\pi_1(X,a)\]
\[\sigma_1 \tensor \cdots \tensor \sigma_q \mapsto (1_a- \sigma_1)\cdot \cdots \cdot (1_a - \sigma_q)\]
descends to an isomorphism on zeroth homology. Then it is clear that the tensor length truncations of $\Cobar$ correspond to the $\mathcal{I}$-adic filtration on the fundamental group ring, since the latter is generated by elements of the form $1-g$. The natural generalization of Stallings' map to the two-sided case 
\[\tot_0(\Cobar(b,X,a)) \to \Z\pi_X(a,b)\]
\[1_b \tensor \sigma_1 \tensor \cdots \tensor \sigma_q \tensor 1_a \mapsto (1- \sigma_1)\cdot \cdots \cdot (1- \sigma_q)\]
runs into a few issues. First, it is unclear what the multiplication on the right hand side means. In the fundamental group ring the multiplication is given by path concatenation, and all the $\sigma_i$ are loops at $a$, but once we have multiple zero-simplicies this becomes ill-defined. Secondly, recall that $\tot_0\Cobar(b,X,a)$ is generated by elements in 
\[(C_\bullet(b) \tensor \underbrace{C_\bullet(X) \tensor \cdots \tensor C_\bullet(X)}_{\text{$q$ times}} \tensor C_\bullet(a))_q,\]
so one could, for example, have elements in $\tot_0(\Cobar(b,X,a))$ from
\[C_0(b) \tensor C_0(X)\tensor C_2(X)\tensor C_0(a),\]
in which case it is unclear how Stallings' map generalizes, since we have only defined it on 1-chains. 

To address the first issue, we make some auxilliary constructions. Recall that the \emph{fundamental groupoid} of a space $X$ is the category $\Pi_X$ whose objects are points in $X$ and morphisms are homotopy classes of paths between two points. 
\begin{defn}
     Let $\Z\Pi^0(X)$ denote the category whose objects are zero-simplicies of $X$ and morphisms are the free abelian group on homotopy classes of paths.
\end{defn}
To be able to define multiplication in this category, we regard it as a sort of ''multiplicative subset'' of a ring as follows. 
\begin{defn}
For any small category $\mathcal{C}$ and ring $R$, the \emph{category algebra} $R\mathcal{C}$ is the ring (possibly without multiplicative identity) whose underlying set is the free $R$-module on the set of all morphisms in $\mathcal{C}$. The addition on $\R\mathcal{C}$ is given by the free module addition, and the multiplication of two morphisms is given by their composition, if they are composable, and zero otherwise. 
\end{defn}
If $\mathcal{C}$ has only finitely many objects, then the sum of all identity objects serves as the multiplicative identity. This also makes $R\mathcal{C}$ into an $R$-algebra. In our case, let $X$ be a topological space, and first consider the full subcategory of the fundamental groupoid whose objects are just the zero simplicies of $X$, which we will denote by $\Pi_X^0$. Form the integral category algebra $\Z\Pi_X^0$ on this category. An element of $\Z\Pi_X^0$ is some $\Z$-linear combination of paths between zero simplicies, and multiplication of two homotopy classes of paths $\sigma$ and $\tau$ is given by concatenation if the endpoint of $\sigma$ is equal to the starting point of $\tau$, and zero otherwise. By regarding $\Z\Pi(X)(a,b)$ as a subset of $\Z\Pi_X^0$, we can multiply its elements. Indeed, in the case with a single zero-simplex, the mutiplication reduces to path concatenation in the fundamental group ring.
\begin{prop}\label{prop:phi_a_b}
    For $a,b \in X^0$, the map 
    \[\varphi_{a,b}: \tot_0\Cobar^{\leq n}(N_\bullet(b), N_\bullet(X), N_\bullet(a)) \to \Z\pi_X(a,b)/\Z\pi_X(a,b)\mathcal{I}^{n+1}\]
    given by
    \[\varphi(1_a \tensor \sigma_1 \tensor \cdots\tensor \sigma_q \tensor 1_q) = 1_a\cdot (\prod_{i=1}^q(1-\sigma_i))\cdot 1_b\]
    on generators consisting of 1-simplicies, and zero otherwise, descends to an isomorphism on zeroth homology.
    \begin{proof}
        Let us first check that $\varphi_{a,b}$ descends to homology. By definition, $\tot_1\Cobar$ necessarily contains a tensor factor coming from $N_2(X)$, say $\Lambda$. Computing the two-sided cobar differential, we have 
        \begin{equation}
            \begin{split}  
            \varphi_{a,b}(\partial \Lambda) &= (1 - d_0\Lambda) - (1 - d_1\Lambda) + (1 - d_2\Lambda) \\
            &= 1 - d_0\Lambda + d_1\Lambda - d_2\Lambda \\
            G'_{a,b}(\overline{\Delta}\Lambda) &= (1 - d_2\Lambda)(1 - d_0\Lambda) \\
            &= 1 - d_0\Lambda - d_2\Lambda + d_2\Lambda * d_0\Lambda\\
            \end{split}
        \end{equation}
        so that 
        \[\varphi_{a,b}\circ \partial_{\Cobar}(1_a\tensor \Lambda \tensor 1_b) = G'(\partial\Lambda-\overline{\Delta}(\Lambda))=[d_1\Lambda] -[d_2\Lambda * d_0\Lambda] = 0\]
        since $d_2\Lambda * d_0\Lambda$ is homotopic to $d_1\Lambda$. Thus $\varphi_{a,b}$ descends to a well defined map on zeroth homology.
            
        
        The remainder of the proof relies on a comparison with Stallings' theorem. First, recall that we have a homotopy equivalence $|\Sing(X)|$ from the geometric realization of its singular complex to $X$. Inside $\Sing(X)$, we can contract a maximal spanning tree of zero-simplices $T$, giving a quasi-isomorphism
        \[q: N_\bullet(X) = N_\bullet(|\Sing(X)|) \to N_\bullet(\Sing(X)/T)\]
        where $N_\bullet$ denotes normalized singular chains, and $\Sing(X)/T$ has a single zero-simplex $*$. Let $D$ denote the two-sided cobar bicomplex 
        \[\Cobar(N_\bullet(b), N_\bullet(X), N_\bullet(a))\] 
        and $D^{>n}$, $D^{\leq n}$ be the previous subcomplex and quotient. Then $q$ induces a map from $D$ onto the two-sided cobar construction on
        \begin{equation}\label{eqn:two_sided_cobar_on_contracted}
        (N_\bullet(*), N_\bullet(\Sing(X)/T), N_\bullet(*))
        \end{equation}
        By Lemma \ref{lem:quasiiso_bicomplex} and Proposition \ref{prop:twoside_reduction}, \ref{eqn:two_sided_cobar_on_contracted} is quasi-isomorphic after totalization to the single-sided cobar construction $\Cobar(C_\bullet(|\Sing(X)/T|))$ with basepoint $*$. Let $E$, $E^{>n}$, $E^{\leq n}$ denote the corresponding bicomplexes for $\Cobar(C_\bullet(|\Sing(X)/T|))$. This gives us a quasi-isomorphism $\tot D^{\leq n}\to E^{\leq n}$. Consider now the following diagram:
        % https://q.uiver.app/#q=WzAsNCxbMCwwLCJIXzAoRF57XFxsZXEgbn0pIl0sWzAsMSwiXFxaXFxwaV9YKGEsYikvXFxaXFxwaV9YKGEsYilcXG1hdGhjYWx7SX1ee24rMX0iXSxbMSwwLCJIXzAoRV57XFxsZXEgbn0pIl0sWzEsMSwiXFxaXFxwaV8xKHxcXFNpbmcoWCkvVHwsICopL1xcbWF0aGNhbHtJfV57bisxfSJdLFswLDEsIlxcdmFycGhpX3thLGJ9IiwyLHsic3R5bGUiOnsiYm9keSI6eyJuYW1lIjoiZGFzaGVkIn19fV0sWzIsMywiXFxyZWZ7dGhtOnN0YWxsaW5nc30iXSxbMCwyXSxbMSwzXV0=
        \[\begin{tikzcd}[cramped]
            {H_0(D^{\leq n})} & {H_0(E^{\leq n})} \\
            {\Z\pi_X(a,b)/\Z\pi_X(a,b)\mathcal{I}^{n+1}} & {\Z\pi_1(|\Sing(X)/T|, *)/\mathcal{I}^{n+1}}
            \arrow[from=1-1, to=1-2]
            \arrow["{\varphi_{a,b}}"', dashed, from=1-1, to=2-1]
            \arrow["{\ref{thm:stallings}}", from=1-2, to=2-2]
            \arrow[from=2-1, to=2-2]
        \end{tikzcd}\]
        By the previous discussion, the top and right arrows are isomorphisms. The bottom isomorhpism induced by the following map $\Z\pi_X(a,b)\to \Z\pi_1(|\Sing(X)/T|,*)$. A generator $\sigma \in \Z\Pi_X(a,b)$ is a 1-simplex, hence a $1$-cell in $|\Sing(X)|$ which can be concatenated with the path from $b\to a$ in the tree $|T|$. Then quotienting yields an element in $\Z\pi_1(|\Sing(X)/T|, *)$. Since $|T|$ is a maximal tree this is an isomorphism. Since $\varphi_{a,b}$ commutes with these maps, we are done.
    \end{proof}
\end{prop}
We assure the reader that the collection of maps $\varphi_{a,b}$ for $a,b \in X^0$ are natural in the category of topological spaces with two basepoints. 
\begin{proof}[Proof of  \ref{thm:stallings_generalization}]
    First note that the short exact sequence $0 \to D^{>n} \to D \to D^{\leq n} \to 0$ induces a long exact sequence 
\[\cdots \to H_0( D^{>n}) \to H_0(D) \to H_0(D^{\leq n}) \to H_{-1}(D^{>n})\to \cdots.\]
To see that the negative homology is zero, let us work with the normalized two-sided cobar construction, which is quasi-isomorphic to $D$. Then the elements of total degree $k$ belong to 
\[(N_\bullet(b) \tensor \overline{N_\bullet(X)}^{\tensor q} \tensor \overline N_\bullet(a))_p\]
where $k=-p-q$.
Because $N_\bullet(b)$ and $N_\bullet(a)$ are concentrated in degree zero, and $\overline{N_0(X)} = 0$ by path-connectedness, any element in negative total degree must be zero. Thus we have a surjection 
\[H_0(D)\to H_0(D^{\leq n})\]
whose kernel is $H_0(D^{>n})$, i.e. the $n$th filtered piece of $H_0(D)$. Thus $H_0(D) / F_nH_0(D) \cong H_0(D^{\leq n})$. By Proposition \ref{prop:phi_a_b}, we conclude 
\begin{equation}
\begin{split}
D^{\infty}_{-n,n} &\cong F_nH_0(D) / F_{n+1}H_0(D) \\
&\cong \ker(H_0(D^{\leq n}) \to H_0(D^{\leq n-1})) \\
&\cong \ker(\Z\pi_X(a,b)/\Z\pi_X(a,b)\mathcal{I}^{n+1} \to \Z\pi_X(a,b)/\Z\pi_X(a,b)\mathcal{I}^{n}) \\
&=\Z\pi_X(a,b)\mathcal{I}^{n} / \Z\pi_X(a,b)\mathcal{I}^{n+1}
\end{split}
\end{equation}
which completes the proof.
\end{proof}